\markdownRendererDocumentBegin
\markdownRendererSectionBegin
\markdownRendererHeadingOne{Longest Common Subsequence \markdownRendererDollarSign{}(LCS)\markdownRendererDollarSign{}:}\markdownRendererInterblockSeparator
{}LCS encuentra la subsecuencia común más larga entre dos secuencias mediante programación dinámica con complejidad \markdownRendererDollarSign{}O(nm)\markdownRendererDollarSign{}. La relación de recurrencia establece que \markdownRendererDollarSign{}LCS\markdownRendererWarning{Undefined link reference "j"}{Undefined link reference "j"}{Look for the text `[...][j]`.}{Look for the text `[...][j]`.}[i][j] = LCS\markdownRendererWarning{Undefined link reference "j-1"}{Undefined link reference "j-1"}{Look for the text `[...][j-1]`.}{Look for the text `[...][j-1]`.}[i-1][j-1] + 1\markdownRendererDollarSign{} si los caracteres coinciden, o \markdownRendererDollarSign{}max(LCS\markdownRendererWarning{Undefined link reference "j"}{Undefined link reference "j"}{Look for the text `[...][j]`.}{Look for the text `[...][j]`.}[i-1][j], LCS\markdownRendererWarning{Undefined link reference "j-1"}{Undefined link reference "j-1"}{Look for the text `[...][j-1]`.}{Look for the text `[...][j-1]`.}[i][j-1])\markdownRendererDollarSign{} en caso contrario. Esta técnica tiene aplicaciones en bioinformática para alineamiento de secuencias, control de versiones y detección de similitud textual.\markdownRendererInterblockSeparator
{}\markdownRendererInputFencedCode{./_markdown_main/a68e51205005627935d04435c60a8e5f.verbatim}{cpp}{cpp}
\markdownRendererSectionEnd \markdownRendererDocumentEnd